\documentclass{article}
\usepackage[utf8]{inputenc}
\usepackage{array}
\usepackage{listings}
\usepackage{xcolor}
\usepackage{booktabs}
\usepackage{geometry}
\geometry{a4paper, margin=1in}

\begin{document}

\section*{Tablas del Modelo Relacional}

% --- Tabla Deporte ---
\subsection*{Tabla: Deporte}
\textit{Almacena los deportes representados en las tarjetas.}
\begin{center}
\begin{tabular}{|l|l|p{6cm}|}
\hline
\textbf{Nombre del Campo} & \textbf{Tipo de Dato (Sugerido)} & \textbf{Descripción / Restricción} \\
\hline
\texttt{id\_deporte} & \texttt{INTEGER} & Llave Primaria (PK), Autoincremental \\
\texttt{nombre\_deporte} & \texttt{VARCHAR(50)} & Nombre único del deporte \\
\texttt{descripcion} & \texttt{TEXT} & Descripción del deporte \\
\hline
\end{tabular}
\end{center}

% --- Tabla Equipo ---
\subsection*{Tabla: Equipo}
\textit{Contiene la información de cada equipo.}
\begin{center}
\begin{tabular}{|l|l|p{6cm}|}
\hline
\textbf{Nombre del Campo} & \textbf{Tipo de Dato (Sugerido)} & \textbf{Descripción / Restricción} \\
\hline
\texttt{id\_equipo} & \texttt{INTEGER} & Llave Primaria (PK), Autoincremental \\

\texttt{nombre\_equipo} & \texttt{VARCHAR(100)} & Nombre único del equipo \\
\texttt{ciudad\_origen} & \texttt{VARCHAR(100)} & Ciudad de origen del equipo \\
\texttt{fecha\_primer\_partido} & \texttt{DATE} & Fecha del primer partido jugado \\
\texttt{fecha\_ultimo\_partido} & \texttt{DATE} & Fecha del último partido jugado \\
\texttt{id\_deporte} & \texttt{INTEGER} & Llave Foránea (FK) hacia Deporte \\
\hline
\end{tabular}
\end{center}

% --- Tabla Puesto ---
\subsection*{Tabla: Puesto}
\textit{Define las posiciones de los jugadores y sus detalles asociados.}
\begin{center}
\begin{tabular}{|l|l|p{6cm}|}
\hline
\textbf{Nombre del Campo} & \textbf{Tipo de Dato (Sugerido)} & \textbf{Descripción / Restricción} \\
\hline
\texttt{id\_puesto} & \texttt{INTEGER} & Llave Primaria (PK), Autoincremental \\
\texttt{nombre\_puesto} & \texttt{VARCHAR(100)} & Nombre de la posición \\
\texttt{direccion\_entrenamiento} & \texttt{VARCHAR(150)} & Dirección del lugar de entrenamiento \\
\texttt{telefono\_particular} & \texttt{VARCHAR(20)} & Teléfono particular de contacto \\
\texttt{telefono\_celular} & \texttt{VARCHAR(20)} & Teléfono celular de contacto \\
\texttt{ciudad\_entrenamiento} & \texttt{VARCHAR(100)} & Ciudad del entrenamiento \\
\texttt{estado\_entrenamiento} & \texttt{VARCHAR(100)} & Estado o región del entrenamiento \\
\hline
\end{tabular}
\end{center}

% --- Tabla Jugador ---
\subsection*{Tabla: Jugador}
\textit{Almacena los datos personales y profesionales de cada jugador.}
\begin{center}
\begin{tabular}{|l|l|p{6cm}|}
\hline
\textbf{Nombre del Campo} & \textbf{Tipo de Dato (Sugerido)} & \textbf{Descripción / Restricción} \\
\hline
\texttt{id\_jugador} & \texttt{INTEGER} & Llave Primaria (PK), Autoincremental \\
\texttt{nombre\_jugador} & \texttt{VARCHAR(100)} & Nombre completo del jugador \\
\texttt{genero} & \texttt{VARCHAR(10)} & Género del jugador \\
\texttt{fecha\_nacimiento} & \texttt{DATE} & Fecha de nacimiento \\
\texttt{lugar\_nacimiento} & \texttt{VARCHAR(150)} & Lugar de nacimiento \\
\texttt{fecha\_pertenencia\_equipo} & \texttt{DATE} & Fecha de ingreso al equipo actual \\
\texttt{id\_equipo} & \texttt{INTEGER} & Llave Foránea (FK) hacia Equipo \\
\texttt{id\_puesto} & \texttt{INTEGER} & Llave Foránea (FK) hacia Puesto \\
\hline
\end{tabular}
\end{center}

% --- Tabla Tarjeta ---
\subsection*{Tabla: Tarjeta}
\textit{Contiene toda la información específica de cada tarjeta coleccionable.}
\begin{center}
\begin{tabular}{|l|l|p{6cm}|}
\hline
\textbf{Nombre del Campo} & \textbf{Tipo de Dato (Sugerido)} & \textbf{Descripción / Restricción} \\
\hline
\texttt{id\_tarjeta} & \texttt{INTEGER} & Llave Primaria (PK), Autoincremental \\
\texttt{numero\_tarjeta} & \texttt{VARCHAR(50)} & Número o código de la tarjeta \\
\texttt{lugar\_entrenamiento} & \texttt{VARCHAR(150)} & Lugar específico mostrado en la tarjeta \\
\texttt{fecha\_salida\_venta} & \texttt{DATE} & Fecha de lanzamiento al mercado \\
\texttt{fecha\_venta} & \texttt{DATE} & Fecha en que fue vendida/adquirida \\
\texttt{medio\_adquisicion} & \texttt{VARCHAR(100)} & Cómo fue adquirida (tienda, cambio) \\
\texttt{descripcion} & \texttt{TEXT} & Descripción o datos en la tarjeta \\
\texttt{estado\_coleccion} & \texttt{VARCHAR(100)} & Condición de la tarjeta (nuevo, usado) \\
\texttt{id\_jugador} & \texttt{INTEGER} & Llave Foránea (FK) hacia Jugador \\
\texttt{id\_equipo} & \texttt{INTEGER} & Llave Foránea (FK) hacia Equipo \\
\hline
\end{tabular}
\end{center}


\vspace{5mm}

\section{Modelo Relacional}

\section{Descripción de las Relaciones}

\begin{enumerate}
    \item \textbf{Deporte - Equipo} (1:N)
    \begin{itemize}
        \item Un deporte puede tener múltiples equipos
        \item Un equipo pertenece a un solo deporte
    \end{itemize}
    
    \item \textbf{Equipo - Jugador} (1:N)
    \begin{itemize}
        \item Un equipo tiene múltiples jugadores
        \item Un jugador pertenece a un solo equipo
    \end{itemize}
    
    \item \textbf{Puesto - Jugador} (1:N)
    \begin{itemize}
        \item Un puesto puede ser ocupado por múltiples jugadores
        \item Un jugador tiene un solo puesto principal
    \end{itemize}
    
    \item \textbf{Jugador - Tarjeta} (1:N)
    \begin{itemize}
        \item Un jugador puede tener múltiples tarjetas
        \item Una tarjeta representa a un solo jugador
    \end{itemize}
    
    \item \textbf{Equipo - Tarjeta} (1:N)
    \begin{itemize}
        \item Un equipo puede estar representado en múltiples tarjetas
        \item Una tarjeta representa a un solo equipo
    \end{itemize}
\end{enumerate}

\section{Atributos Clave}

\begin{itemize}
    \item \textbf{Deporte}: id\_deporte (PK)
    \item \textbf{Equipo}: id\_equipo (PK), id\_deporte (FK)
    \item \textbf{Puesto}: id\_puesto (PK)
    \item \textbf{Jugador}: id\_jugador (PK), id\_equipo (FK), id\_puesto (FK)
    \item \textbf{Tarjeta}: id\_tarjeta (PK), id\_jugador (FK), id\_equipo (FK)
\end{itemize}

\section{Restricciones Identificadas}

\begin{enumerate}
    \item Unicidad en nombre\_deporte
    \item Unicidad en nombre\_equipo por deporte
    \item Fechas válidas (fecha\_ultimo\_partido $\geq$ fecha\_primer\_partido)
    \item Formato de teléfonos
    \item Valores específicos para género y estado\_coleccion
    \item Integridad referencial entre todas las tablas relacionadas
\end{enumerate}





\vspace{5mm}

\section{Modelo Relacional Normalizado}
\section*{Modelo Relacional Normalizado}

% --- Tabla Deporte ---
\subsection*{Tabla: Deporte}
\begin{center}
\begin{tabular}{|l|l|p{6cm}|}
\hline
\textbf{Campo} & \textbf{Tipo de Dato} & \textbf{Restricción} \\
\hline
\texttt{id\_deporte} & INTEGER & PK, Autoincremental \\
\texttt{nombre\_deporte} & VARCHAR(50) & NOT NULL, Único \\
\texttt{descripcion} & TEXT &  \\
\hline
\end{tabular}
\end{center}

% --- Tabla Equipo ---
\subsection*{Tabla: Equipo}
\begin{center}
\begin{tabular}{|l|l|p{6cm}|}
\hline
\textbf{Campo} & \textbf{Tipo de Dato} & \textbf{Restricción} \\
\hline
\texttt{id\_equipo} & INTEGER & PK, Autoincremental \\
\texttt{nombre\_equipo} & VARCHAR(100) & NOT NULL, Único \\
\texttt{ciudad\_origen} & VARCHAR(100) &  \\
\texttt{fecha\_primer\_partido} & DATE &  \\
\texttt{fecha\_ultimo\_partido} & DATE &  \\
\texttt{id\_deporte} & INTEGER & FK $\rightarrow$ Deporte(id\_deporte) \\
\hline
\end{tabular}
\end{center}

% --- Tabla Puesto ---
\subsection*{Tabla: Puesto}
\begin{center}
\begin{tabular}{|l|l|p{6cm}|}
\hline
\textbf{Campo} & \textbf{Tipo de Dato} & \textbf{Restricción} \\
\hline
\texttt{id\_puesto} & INTEGER & PK, Autoincremental \\
\texttt{nombre\_puesto} & VARCHAR(100) & NOT NULL \\
\texttt{direccion\_entrenamiento} & VARCHAR(150) &  \\
\texttt{telefono\_particular} & number(1) &  \\
\texttt{telefono\_celular} & VARCHAR(20) &  \\
\texttt{ciudad\_entrenamiento} & VARCHAR(100) &  \\
\texttt{estado\_entrenamiento} & VARCHAR(100) &  \\
\hline
\end{tabular}
\end{center}

% --- Tabla Jugador ---
\subsection*{Tabla: Jugador}
\begin{center}
\begin{tabular}{|l|l|p{6cm}|}
\hline
\textbf{Campo} & \textbf{Tipo de Dato} & \textbf{Restricción} \\
\hline
\texttt{id\_jugador} & INTEGER & PK, Autoincremental \\
\texttt{nombre\_jugador} & VARCHAR(100) & NOT NULL \\
\texttt{genero} & VARCHAR(10) &  \\
\texttt{fecha\_nacimiento} & DATE &  \\
\texttt{lugar\_nacimiento} & VARCHAR(150) &  \\
\texttt{fecha\_pertenencia\_equipo} & DATE &  \\
\texttt{id\_equipo} & INTEGER & FK $\rightarrow$ Equipo(id\_equipo) \\
\texttt{id\_puesto} & INTEGER & FK $\rightarrow$ Puesto(id\_puesto) \\
\hline
\end{tabular}
\end{center}

% --- Tabla Tarjeta ---
\subsection*{Tabla: Tarjeta}
\begin{center}
\begin{tabular}{|l|l|p{6cm}|}
\hline
\textbf{Campo} & \textbf{Tipo de Dato} & \textbf{Restricción} \\
\hline
\texttt{id\_tarjeta} & INTEGER & PK, Autoincremental \\
\texttt{numero\_tarjeta} & VARCHAR(50) & NOT NULL, Único \\
\texttt{lugar\_entrenamiento} & VARCHAR(150) &  \\
\texttt{fecha\_salida\_venta} & DATE &  \\
\texttt{fecha\_venta} & DATE &  \\
\texttt{medio\_adquisicion} & VARCHAR(100) &  \\
\texttt{descripcion} & TEXT &  \\
\texttt{estado\_coleccion} & VARCHAR(100) &  \\
\texttt{id\_jugador} & INTEGER & FK $\rightarrow$ Jugador(id\_jugador) \\
\texttt{id\_equipo} & INTEGER & FK $\rightarrow$ Equipo(id\_equipo) \\
\hline
\end{tabular}
\end{center}






\lstset{
    language=SQL,
    basicstyle=\ttfamily\small,
    keywordstyle=\color{blue}\bfseries,
    commentstyle=\color{gray}\itshape,
    stringstyle=\color{red},
    showstringspaces=false,
    breaklines=true,
    frame=single,
    tabsize=2
}

\section*{Listados de Programas en SQL}

\begin{itemize}
    \item Creación de base de datos con el nombre que se indica en el encabezado.
    \item Creación de tablas.
    \item Creación de llaves foráneas y primarias según corresponda con el modelo lógico.
    \item Creación de reglas, constraints o defaults que permitan integridad de dominio y/o reglas de negocio.
    \item Alta de cuando menos 10 registros a las tablas catálogo de la BD y los registros correspondientes en las tablas relación para denotar las asociaciones 1:1; 1:N; M:N. Donde M y N son cuando menos 3.
    \item Creación de índices en campos líderes de consulta común.
    \item Procedimiento almacenado \texttt{sp\_inserta} que, a través de parámetros permita inserción de registros de alguna tabla catálogo.
    \item Creación de consultas, que satisfagan la hoja de requerimientos.
    \item Procedimiento almacenado \texttt{sp\_consultas} que acepte parámetros y ejecute 5 consultas, si los requerimientos contienen menos crear dos nuevas. La salida del Procedimiento (es decir, el resultado de las consultas) debe guardarlo en un archivo.
    \item Procedimiento almacenado \texttt{sp\_borra} que permita eliminación de registros de alguna tabla catálogo.
    \item Comandos necesarios para el control de la integridad de la información.
\end{itemize}

\subsection*{Código SQL}

\begin{lstlisting}[language=SQL]
-- Creación de la base de datos
CREATE DATABASE Colecciones;
GO

USE Colecciones;
GO

-- Tabla Deporte
CREATE TABLE Deporte (
    id_deporte INT PRIMARY KEY IDENTITY(1,1),
    nombre_deporte VARCHAR(50) UNIQUE NOT NULL,
    descripcion TEXT
);

-- Tabla Equipo
CREATE TABLE Equipo (
    id_equipo INT PRIMARY KEY IDENTITY(1,1),
    nombre_equipo VARCHAR(100) UNIQUE NOT NULL,
    ciudad_origen VARCHAR(100),
    fecha_primer_partido DATE,
    fecha_ultimo_partido DATE,
    id_deporte INT NOT NULL,
    FOREIGN KEY (id_deporte) REFERENCES Deporte(id_deporte)
);

-- Tabla Puesto
CREATE TABLE Puesto (
    id_puesto INT PRIMARY KEY IDENTITY(1,1),
    nombre_puesto VARCHAR(100) NOT NULL,
    direccion_entrenamiento VARCHAR(255),
    telefono_particular VARCHAR(20),
    telefono_celular VARCHAR(20),
    ciudad_entrenamiento VARCHAR(100),
    estado_entrenamiento VARCHAR(100)
);

-- Tabla Jugador
CREATE TABLE Jugador (
    id_jugador INT PRIMARY KEY IDENTITY(1,1),
    nombre_jugador VARCHAR(100) NOT NULL,
    genero VARCHAR(10),
    fecha_nacimiento DATE,
    lugar_nacimiento VARCHAR(150),
    fecha_pertenencia_equipo DATE,
    id_equipo INT NOT NULL,
    id_puesto INT NOT NULL,
    FOREIGN KEY (id_equipo) REFERENCES Equipo(id_equipo),
    FOREIGN KEY (id_puesto) REFERENCES Puesto(id_puesto)
);

-- Tabla Tarjeta
CREATE TABLE Tarjeta (
    id_tarjeta INT PRIMARY KEY IDENTITY(1,1),
    numero_tarjeta VARCHAR(50) UNIQUE NOT NULL,
    lugar_entrenamiento VARCHAR(255),
    fecha_salida_venta DATE,
    fecha_venta DATE,
    medio_adquisicion VARCHAR(100),
    descripcion TEXT,
    estado_coleccion VARCHAR(100),
    id_jugador INT NOT NULL,
    id_equipo INT NOT NULL,
    FOREIGN KEY (id_jugador) REFERENCES Jugador(id_jugador),
    FOREIGN KEY (id_equipo) REFERENCES Equipo(id_equipo)
);

-- Ejemplo de inserciones (mínimo 10 registros en catálogos)
INSERT INTO Deporte (nombre_deporte, descripcion) VALUES
('Futbol', 'Deporte en el que se enfrentan dos equipos de 11 jugadores'),
('Baseball', 'Deporte con bate y pelota'),
('Basketball', 'Deporte con balón y canasta');

-- Ejemplo de índice
CREATE INDEX idx_nombre_equipo ON Equipo(nombre_equipo);

-- Procedimiento almacenado para inserción
CREATE PROCEDURE sp_inserta_equipo
    @nombre_equipo VARCHAR(100),
    @ciudad_origen VARCHAR(100),
    @id_deporte INT
AS
BEGIN
    INSERT INTO Equipo (nombre_equipo, ciudad_origen, id_deporte)
    VALUES (@nombre_equipo, @ciudad_origen, @id_deporte);
END;
GO

-- Procedimiento almacenado para consultas
CREATE PROCEDURE sp_consultas
    @opcion INT
AS
BEGIN
    IF @opcion = 1
        SELECT * FROM Deporte;
    ELSE IF @opcion = 2
        SELECT * FROM Equipo;
    ELSE IF @opcion = 3
        SELECT * FROM Jugador;
    ELSE IF @opcion = 4
        SELECT * FROM Tarjeta;
    ELSE IF @opcion = 5
        SELECT d.nombre_deporte, COUNT(e.id_equipo) AS total_equipos
        FROM Deporte d
        JOIN Equipo e ON d.id_deporte = e.id_deporte
        GROUP BY d.nombre_deporte;
END;
GO

-- Procedimiento almacenado para borrado
CREATE PROCEDURE sp_borra_deporte
    @id_deporte INT
AS
BEGIN
    DELETE FROM Deporte WHERE id_deporte = @id_deporte;
END;
GO
\end{lstlisting}







\end{document}